\documentclass[11pt]{article}
\usepackage[a4paper,margin=26mm]{geometry}
\usepackage{amsmath,amssymb,amsthm,mathtools,booktabs,hyperref,xurl}
\setlength{\emergencystretch}{3em}
\newtheorem{theorem}{Theorem}
\newtheorem{lemma}[theorem]{Lemma}
\newtheorem{proposition}[theorem]{Proposition}
\newtheorem{corollary}[theorem]{Corollary}
\newtheorem{remark}[theorem]{Remark}
\newcommand{\Q}{\mathbb Q}
\newcommand{\Z}{\mathbb Z}
\newcommand{\N}{\mathbb N}
\newcommand{\Tr}{\operatorname{Tr}}
\newcommand{\Nm}{\operatorname{N}}
\title{A common numerator point for simultaneous integer certificates\\
and rational undecidability in 83, 101, and 53 variables}
\author{\(\mathbb Z\mathbb Q\) Diophantine Project}
\date{23 November 2026}
\begin{document}
\maketitle
\begin{abstract}
We improve the simultaneous integer certificates of the preceding
99-, 117-, and 69-variable transfer by using one numerator point for
all input blocks. The numerator index is selected independently of the
reference index. Disjoint contact families and an all-sign coprimality
property of the retained prime supplier allow the same numerator point
to approximate different integer multiples of the reference point at
different selected primes. This saves two rational coordinates for
every input after the first.

The resulting joint certificates for \(\Z^k\) use \(7k+18\) coordinates
in the general two-detector route, \(9k+18\) in the specified
three-detector route, and \(5k+6\) in the ordinary-\(abc\) route.
The mixed-integer universality theorem gives 83, 101, and 53 rational
unknowns, with degree bounds strictly below \(2^{413},2^{477}\), and
\(2^{263}\), respectively. All retained non-pole arithmetic premises
are stated explicitly; the first two routes introduce no additional
\(abc\) hypothesis. These are general-degree statements, not quartic
undecidability bounds. Complete coefficient circuits, source maps,
guard lists, and exact degree certificates accompany the proof.
\end{abstract}

\section{Statement and retained premises}

We use the project sources with full auxiliary elliptic charts,
the four-contact scalar or twelve-contact construction of v78,
and the normalized Bhargava--Burungale product poles of v88
\cite{v78,v88}. The preceding simultaneous construction is
\cite{shared}. Its proof and complete coefficient circuits are
retained. The new operation identifies its private numerator point
across input blocks. We prove that this identification preserves
completeness; soundness will follow by projection to each unchanged
source block.

The arithmetic dependencies retained from these sources are the
elliptic-group and chart interpretations, the stated trace or effective
abscissa-sum height and contact-height bounds, the finite prime-pattern
supplier, and the compatible finite-place and compact-real recurrence
statements. The ordinary-\(abc\) route also retains ordinary rational
\(abc\). The prime supplier is used for every fixed finite rational
contact list, with arbitrary fixed divisibility and finite avoidance;
its all-sign order-coprimality assertion is given explicitly in
Section~\ref{sec:crt}. These arithmetic dependencies are not proved
by the coefficient verification.

The v88 product-pole completion uses the published
Bhargava--Burungale theorem \cite{BB}. It does not retain the earlier
native-pole realization premise. All pole parameters remain private
to their original input and detector. In particular, this proof
does not identify distinct square classes or request a common pole
witness.

\begin{theorem}\label{thm:main}
Under the corresponding retained source premises, for every fixed
integer \(k\ge1\) there is an integer polynomial with \(k\)
distinguished rational inputs whose input projection is exactly
\(\Z^k\), with the following total numbers of rational coordinates:
\[
\begin{array}{c|c|c}
\text{route}&\text{coordinates}&\text{source parameter}\\ \hline
\text{general two-detector}&7k+18&\text{sufficiently large fixed odd }L\\
\text{specified three-detector}&9k+18&L=143\\
\text{ordinary }abc&5k+6&\text{sufficiently large fixed odd }B .
\end{array}
\]
The sufficient thresholds in the first and third lines are the
single-input soundness thresholds; they are not increased with \(k\).
Consequently rational polynomial solvability is undecidable in
83, 101, and 53 unknowns, respectively, under the stated premises.
The final polynomials may be taken to have degrees strictly below
\(2^{413},2^{477}\), and \(2^{263}\), respectively.
\end{theorem}

The first and third assertions are existence statements for a fixed
defining member when the sufficient parameter threshold is ineffective.
The stored general fixture \(L=143\) and ordinary-\(abc\) fixture \(B=3\)
do not certify sufficient parameters for those two routes. The specified
three-detector \(L=143\) member is the one certified by its source.
All degree bounds are uniform in these fixed coefficient parameters.
These qualifications are part of the theorem, not exceptions supplied
by the finite verifier.

\section{Point roles and the new choice}
\label{sec:roles}

Put
\[
 K=\Q(\sqrt2),\qquad D=11-6\sqrt2,\qquad E:y^2=x^3-Dx,
\]
\[
 P=\left(\frac{57-28\sqrt2}{8},\frac{-660+417\sqrt2}{32}\right).
\]
The source fixes the interpretation of \(P\) as its generator of
\(2E(K)\). Write
\[
              f(Q)=\frac1{x(Q)},\qquad z(Q)=-\frac{x(Q)}{y(Q)}.
\]
Every division below retains the source denominator guard.
The nonzero integer label for an input \(X_i\in\Z\) is
\[
                    n_i=X_i^2+1,\qquad A_i=[n_i]P .
\]
The root point \(A_i\) remains private. This also handles \(X_i=0\),
for which \(n_i=1\).

In the general routes the common root labels \(a,t,b\) and the
common point roles from \cite{shared} are
\[
\begin{split}
 S&=[a]P,\quad T=[t]P,\quad Z=[b]P,\quad H=[h]P,\\
 U_S&=[ha]P,\quad U_T=[ht]P,\quad U_Z=[hb]P.
\end{split}
\]
Previously the remaining numerator was \(U_{A_i}=[hn_i]P\).
We replace all these numerator roles by a single point
\[
                              U=[j]P,
\]
where \(j\) is a new integer \emph{label in the completeness proof},
not an additional quantified coordinate. It need not equal any
of the integers \(hn_i\). The equations retain the same rational
point chart for the numerator and every original comparison test.
The two input-dependent ratios now are
\begin{equation}\label{eq:ratios}
 r=\frac{z(U)}{z(H)},\qquad
 r_i'=\overline{\frac{z(U-A_i)}{z(H-P)}}.
\end{equation}
The bar denotes conjugation in \(K/\Q\).

The ordinary-\(abc\) source has only the roles \(A_i,H,U_{A_i}\).
After the same identification it has private \(A_i\) and common
\(H,U\). Its original nonzero source input is replaced throughout
by \(n_i=X_i^2+1\); the separate zero-input polynomial factor is
omitted. This is legitimate because
\[
              X\in\Z\quad\Longleftrightarrow\quad X^2+1\in\Z
              \qquad(X\in\Q).
\]
For the reverse implication, write \(X=r/s\) in lowest terms;
integrality forces \(s^2\mid r^2\), and hence \(s=1\).

Each point role has two retained rational chart coordinates.
Identifying the numerator therefore saves \(2(k-1)\) retained
coordinates. Its three finite chart generators are identified as
well, but those generators are eliminated by the finite-algebra
encoder and are not rational unknowns. The remainder of the proof
justifies this saving.

\section{Contact families remain disjoint}
\label{sec:contacts}

We retain exactly the contact substitutions of \cite{shared}.
For clarity we recall why contacts belonging to different inputs
cannot coincide.

In the general routes let \(e=2\) or \(3\) be the number of
detectors. The v88 scalar for block \(i\) is
\[
 \beta_i=1632X_i^2+2477,\qquad
 A_i^{\rm sc}=\beta_i(\beta_i-1),\qquad
 B_i^{\rm sc}=(\beta_i-2401)(32\beta_i+343).
\]
\[
 T_L^{\rm sc}=\bigl(1224(67^{(L+1)/2}+1)\bigr)^2,\qquad
 \rho_i=102^2\,67^L
 \frac{7T_L^{\rm sc}A_i^{\rm sc}B_i^{\rm sc}}
 {(B_i^{\rm sc}+7T_L^{\rm sc}A_i^{\rm sc})^2}.
\]
For integer inputs these are fixed positive numbers with
\(0<\rho_i<1/167\), square reduced denominator, and the source
class-seven and finite-depth properties. Put
\[
 x(A_i)+x(S)+x(T)+x(Z)=\mathcal U_i+\sqrt2\mathcal S_i.
\]
The two-detector scales are
\[
           \alpha_{i0}=1,\qquad
           \alpha_{i1}=1+\rho_i(2\mathcal U_i)^2 ,
\]
and the three-detector scales are
\[
 \alpha_{i0}=1,\quad
 \alpha_{i1}=1+\rho_i\mathcal U_i^2,\quad
 \alpha_{i2}=1+\rho_i(\mathcal U_i+\mathcal S_i)^2 .
\]
All are at least one and depend only on root points and inputs;
in particular they do not depend on the new \(U\).
For \(0\le i<k\) retain the common-parameter substitution
\[
 t_{i\ell}=\alpha_{i\ell}(\lambda+100i),\qquad
 \delta_{i\ell}=t_{i\ell}+\ell+1 .
\]
With
\[
 \Theta=\{0,-1/2,-1,1/2399,-32/407\},\qquad
 \Theta_\ell=(4\ell+2+\Theta)\cup(4\ell+4+\Theta),
\]
the contact function is
\begin{equation}\label{eq:R}
 R_{i\ell}=
 \frac{\prod_{\theta\in\Theta_\ell}(t_{i\ell}-\theta)}
      {\delta_{i\ell}^{10}}.
\end{equation}
As a function of \(\lambda\), its contact positions are
\[
             d=\theta/\alpha_{i\ell}-100i.
\]
They lie in the pairwise disjoint intervals
\((-100i,-100i+13)\). Repeated positions within one input block
may be merged, keeping all their detector incidences. They impose
the same integer label \(n_i\), so no conflict results.

For the ordinary-\(abc\) source put
\[
 \beta_i=n_i^2+B,\quad T_i=\Tr x(A_i),\quad
 a_i=\beta_i^4+T_i^4,\quad z_i=a_i(z+100i).
\]
Here \(a_i\ge81\) for the source range \(B\ge3\), and \(a_i\) is
fixed before the reference and numerator choices. If
\[
 \mathcal C=\{-9,-8,-7,-5,-4,-3,-2,-1,6,10,11,12\},
\]
the normalized positions are
\[
                      d=\pm(1+c/a_i)-100i,\qquad c\in\mathcal C.
\]
They lie in \((-100i-2,-100i+2)\), again disjoint for different
inputs. No size assumption on \(n_i\) is needed.

Let \(\mathcal D_i\) be the resulting set of distinct positions
assigned to input \(i\), and \(\mathcal D=\bigcup_i\mathcal D_i\).
All these finite sets are fixed before prime selection.

\section{The common-index Chinese remainder theorem}
\label{sec:crt}

The key arithmetic input is the all-sign conclusion of the pinned
Sun prime-pattern source \cite[Theorem 4.1, label
\texttt{an:thm:patterns}]{patterns}. It states the following.
For any distinct integral coefficients \(c_d\), any prescribed
divisibility of \(w\), any odd squarefree avoidance integer \(M_1\)
with \(7\nmid M_1\), and arbitrary positive lower bounds, there
are \(u,w>0\) such that all
\[
                    p_d=u+c_dw^2
\]
are large primes, \(u/w^2\) is arbitrarily large, and their
Gaussian generators satisfy
\begin{equation}\label{eq:gaussian}
 \pi_d\equiv9\pmod {16},\qquad
 \pi_d\equiv1+3\mathrm i\pmod7,\qquad
 \pi_d\equiv\mathrm i\pmod {M_1}.
\end{equation}
The two orders at one position have greatest common divisor four.
At distinct positions any common divisor of either order at one
position and either order at the other is a power of two.
Both orders are prime to \(M_1\). The rational contact formulation
used in the project follows after a fixed denominator dilation.
We use this assertion as part of the retained prime supplier;
we do not add a stronger prime selection hypothesis.

Write the divided annihilator orders as
\[
             u_d=\frac{\Nm(\pi_d-1)}4,\qquad
             v_d=\frac{\Nm(\pi_d+1)}4 .
\]
The symbol \(\mathrm i\) in \eqref{eq:gaussian} is the Gaussian
unit and is unrelated to the input index \(i\).
The normalization in the elliptic source makes \(u_d\) an
annihilator for \(P\) in the primary slot and \(v_d\) an
annihilator in the opposite slot. The residue modulo sixteen gives
\begin{equation}\label{eq:twodepth}
                     v_2(u_d)=4,\qquad v_2(v_d)=0 .
\end{equation}
Indeed, writing \(\pi=a+b\mathrm i\), one has \(a\equiv9\) and
\(b\equiv0\pmod {16}\); hence
\(v_2((a-1)^2+b^2)=6\) and
\(v_2((a+1)^2+b^2)=2\).
It follows that
\begin{equation}\label{eq:ordercoprime}
 \gcd(u_d,v_e)=1\quad(d,e\in\mathcal D),\qquad
 \gcd(v_d,v_e)=1\quad(d\ne e).
\end{equation}
The second assertion uses the same-sign part of the pinned theorem,
not merely its explicitly displayed cross-sign corollary.

\begin{lemma}[Compatibility with a fixed preparation modulus]
\label{lem:avoidance}
Let \(M\) be any fixed positive integer divisible by \(64\).
The prime selection can ensure that the odd parts of every
\(u_d,v_d\) are prime to \(M\), and that all \(p_d\) are prime
to \(M\).
\end{lemma}
\begin{proof}
In \eqref{eq:gaussian} take \(M_1\) to be the product of the
odd prime divisors of \(M\), omitting \(7\).
For each such odd prime \(q\ne7\), the two norms are \(2\pmod q\).
At seven the separate residue in \eqref{eq:gaussian} gives
\[
     \Nm(\pi_d-1)\equiv2\pmod7,\qquad
     \Nm(\pi_d+1)\equiv6\pmod7.
\]
Thus seven is automatically avoided; the restriction
\(7\nmid M_1\) loses nothing. Division by four is harmless at
every odd prime. Finally choose the selected primes larger
than every prime divisor of \(M\).
\end{proof}

We also need a different coprimality fact: every selected prime
must avoid every selected annihilator. Choose a denominator
\(D_0\) for the contact positions, put \(v=D_0^2w^2\), and use
\[
                 \lambda=u/v,\qquad p_d=u-dv
\]
(or the analogous common variable \(z=u/v\)).
The supplier applies to the resulting integral quadratic coefficients.
Impose \(\lambda>4\max_d|d|\) and \(u\ge100\).
Then every selected prime belongs to \((3u/4,5u/4)\).
The source reduction orders are even, and Hasse's bound gives
\[
 \#E(\mathbb F_{p_e})
 \le p_e+1+2\sqrt{p_e}
 <\frac54u+1+\sqrt{5u}<\frac32u.
\]
If \(p_d\) divided one of these even orders, that order would
be at least \(2p_d>3u/2\), a contradiction. Consequently
\begin{equation}\label{eq:primeorder}
                  \gcd(p_d,\prod_eu_ev_e)=1 .
\end{equation}
This cluster argument permits prime-power refinements. It does
not replace the order-order coprimality in
\eqref{eq:ordercoprime}.

\begin{lemma}[Two common indices]\label{lem:indices}
Suppose that \(M\) preserves the fixed-place preparations and
is divisible by \(64\) and by the component-group order of the
compact real cyclic closure used in the source. Fix allowed
initial indices \(h_0,j_0\) divisible by \(64\).
There are positive integers \(h,j\), with independent real
returns to their prescribed allowed neighborhoods, such that
for \(d\in\mathcal D_i\)
\begin{align}
 h&\equiv h_0\pmod M,& h&\equiv0\pmod {u_d},&
 h&\equiv1\pmod {v_d},&
 h&\equiv p_d\pmod {p_d^2},\label{eq:hcrt}\\
 j&\equiv j_0\pmod M,& j&\equiv0\pmod {u_d},&
 j&\equiv n_i\pmod {v_d},&
 j&\equiv hn_i\pmod {p_d^5}.\label{eq:jcrt}
\end{align}
The exponent \(5\) can be replaced by any prescribed larger
finite exponent, separately at each selected prime.
\end{lemma}
\begin{proof}
First apply the generalized Chinese remainder theorem to
\eqref{eq:hcrt}. The only overlap with \(M\) of the primary
orders has two-part \(16\), on which \(h_0\) vanishes. The
primary residues are all zero. The opposite orders are odd,
pairwise coprime, and coprime to all primary orders by
\eqref{eq:ordercoprime}. They are prime to \(M\) by
Lemma~\ref{lem:avoidance}. The selected prime squares are
mutually coprime and prime to every other modulus by
\eqref{eq:primeorder} and finite avoidance. Thus all residues
are compatible.

After selecting \(h\), apply the same argument to
\eqref{eq:jcrt}. The crucial additional fact is that distinct
opposite orders are pairwise coprime, so they may receive
unrelated residues \(n_i\). A repeated contact position belongs
to only one input block and receives only that block's residue.
The residue \(j_0\) agrees with all zero primary residues on
their common two-part. The prime-power conditions are independent
of these order conditions. The same argument works for any
higher prescribed precision.

It remains to preserve the real returns. Let \(G\) be the compact
closure of the cyclic orbit in the product of the source real
groups. If an integer \(m\) is divisible by the order of
\(G/G^0\), then \([m]G=G^0\): multiplication by \(m\) is
surjective on the connected compact torus \(G^0\) and kills
the component group. Hence the orbit of any arithmetic
progression with such step is dense in the component specified
by its residue. The final progressions in \eqref{eq:hcrt} and
\eqref{eq:jcrt} refine the original classes modulo \(M\), so
they retain the prescribed components and admit arbitrarily
large positive indices in the prescribed nonempty open
neighborhoods. The indices are chosen independently.
\end{proof}

\section{Formal estimates at a selected contact}
\label{sec:formal}

Fix \(d\in\mathcal D_i\) and write \(p=p_d\).
Enlarge the fixed avoidance set before prime selection to
contain all bad primes, the supports of every nonzero root
label \(n_i,a,t,b\), all scalars and scales, all contact
differences, \(D_0\), and the preparation modulus and chart
denominators. Thus \(n_i\) is a \(p\)-adic unit.

At the primary slot the congruences modulo \(u_d\) put
\(H=[h]P\) and \(U=[j]P\) in the formal group. At the opposite
slot the congruences modulo \(v_d\) put \(H-P\) and \(U-A_i\)
in the formal group. Set \(\epsilon_i=j-hn_i\). Then
\[
       v_p(h)=1,\qquad v_p(h-1)=0,\qquad
       v_p(\epsilon_i)\ge5.
\]
The exact formal-logarithm ratios are
\begin{equation}\label{eq:exactratios}
 \frac jh=n_i+\frac{\epsilon_i}{h},\qquad
 \frac{j-n_i}{h-1}=n_i+\frac{\epsilon_i}{h-1}.
\end{equation}
This follows by applying the formal logarithm after multiplying
\(P\) by the appropriate annihilator; that annihilator is
prime to \(p\). In particular the first ratio is \(n_i\)
modulo \(p^4\), and the second is \(n_i\) modulo \(p^5\).
Their difference is exactly
\begin{equation}\label{eq:difference}
 \frac jh-\frac{j-n_i}{h-1}
             =-\frac{\epsilon_i}{h(h-1)}
             \equiv0\pmod {p^4}.
\end{equation}

For the source curve, the local parameter satisfies
\(\log_E(z)=z+O(z^5)\) at every retained good odd place.
Thus the ratio of two local parameters is its logarithm ratio
times a unit congruent to one modulo the fourth power of the
smaller parameter ideal. The primary parameters have depth
at least two: since \(p\nmid u_d\) and \(u_d\mid h\), the
point \([h/p]P\) is already formal, and multiplication by \(p\)
adds one to its positive depth. The primary \(U\) has the
same depth because \(j/h\) is a unit. The opposite parameters
have depth at least one, and their logarithm ratio is also a
unit. It follows from \eqref{eq:exactratios}--\eqref{eq:difference}
that the actual ratios in \eqref{eq:ratios} satisfy, in the
successful comparison embedding,
\begin{equation}\label{eq:truncated}
                   r-n_i\in p^4\mathcal O,\qquad
                   r-r_i'\in p^4\mathcal O .
\end{equation}
Here the opposite value is transported by the source conjugation
to the comparison slot, as in its original paired-ratio test.

The primary reciprocal abscissas have depth at least four and
the opposite reciprocal abscissas depth at least two, since
\(f=z^2+O(z^6)\). These are at least the original required
depths. Every other numerator in the general source retains
its exact definition \([ha]P,[ht]P,[hb]P\); their paired
comparisons and Lucas residual estimates are unchanged.
The new numerator therefore satisfies all the same positive
detector valuations at this contact as the old point
\([hn_i]P\).

The reference condition \(h\equiv p\pmod {p^2}\) also retains
the source scalar normalization: the primary reference is
of depth at least two, while \(H\) reduces to \(P\) in the
opposite slot. The six scalar components are functions of
\(H,S\) alone. Their nonvanishing denominator and forbidden
valuation checks are consequently exactly those already proved
for the simultaneous source in \cite{shared}.

Thus every contact of every block has its original successful-slot
inverse detector congruent to one, with the opposite contribution
zero or one. In particular the reciprocal trace remains a unit.
This assertion uses the original positive detector estimates;
it does not assume equality \(U=[hn_i]P\) away from the contact.

\section{Fixed-place and real preparation}
\label{sec:fixed}

One must also check the places prepared \emph{before} the contacts
are selected. A common ratio cannot be arbitrarily close to
distinct \(n_i\) at one fixed place. The detector formulas
do not require such a condition.

\subsection{General two- and three-detector routes}

For fixed integer inputs, the joint scalar preparation in
\cite{shared} first chooses a common initial \(H=[h_0]P\)
deep in the formal groups at the finite union of all fixed
preparation places and near the identity at the real embeddings.
It then chooses \(c\in6\Z_{>0}\), with \(a=c+1\), so that
\([c]P=S-P\) is still closer to the identity. The six rational
components of
\[
                       f(H),\quad f(S-P),\quad z(S-P)/z(H)
\]
meet the finitely many required tolerances simultaneously.
The Lucas recurrence gives fixed positive odd \(t,b\), and
all root points and contact scales are then fixed.

This scalar preparation is independent of the private numerator
point. Choose an allowed initial \(U=[j_0]P\), with \(64\mid j_0\),
avoiding the finitely many poles and excluded point-chart fibers
of \(U\) and every \(U-A_i\). For fixed nonzero labels, the
preimages of forbidden points under multiplication or translation
are finite. The allowed orbit and its real recurrence therefore
permit this choice. Shrink the reference and numerator
neighborhoods so that every raw comparison function of every
block is finite there. On smaller compact real neighborhoods
and sufficiently small finite-place neighborhoods these
finitely many functions are uniformly bounded. Include the
corresponding local congruences in \(M\).

The source's first ten detector factors have the form
\[
             1+\left(\frac{C}{\delta_{i\ell}^2q_{i\ell}}\right)^2,
             \qquad q_{i\ell}=\rho_iR_{i\ell};
\]
its remaining comparison factors have the form
\[
       1+\left(\frac{C^2}
        {\delta_{i\ell}^2q_{i\ell}^{\,2a_r-1}}\right)^2,
                 \qquad a_r\in\{1,3\}.
\]
At each previously prepared finite or real place, let the common
contact parameter tend to infinity. Then
\[
             R_{i\ell}\longrightarrow1,\qquad
             \delta_{i\ell}^{-1}\longrightarrow0,
\]
whereas \(\rho_i\ne0\) is fixed. Uniform boundedness of \(C\)
therefore makes every comparison factor tend to one. This
applies in particular to the input test \(C=r-n_i\), irrespective
of whether its fixed-place valuation is positive, zero, or
negative.

The scalar factors do not have this slope suppression, but their
six components are independent of \(U\) and were already
prepared. They retain precisely the old finite and real bounds.
Thus the common numerator introduces no new fixed-place
approximation requirement on \(r\).

Prescribed divisibility of the prime-pattern denominator \(w\)
makes the contact parameter sufficiently large at every prepared
finite place, simultaneously. The finite avoidance set includes
the supports of the scalars, scales, and common contact
denominator. Fresh prime values force the numerator \(u\) to
be a unit at those places, so increasing the prescribed
denominator valuation has the intended effect. An arbitrarily
large real slope supplies the real limits. The subsequent
refinements \eqref{eq:hcrt}--\eqref{eq:jcrt} preserve the
chosen bounded neighborhoods by their residues modulo \(M\)
and independent real returns.

\subsection{Ordinary rational \(abc\)}

In this route the raw tests are
\[
 f(H),\ f(U),\ \overline{f(H-P)},\
 \overline{f(U-A_i)},\ r-r_i',\ r-n_i .
\]
Let \(C_{i1},\ldots,C_{i4}\) be the first four tests, and
let \(C_{i5},C_{i6}\) be the squares of the last two.
Put
\[
 G_0(t)=\prod_{c\in\mathcal C}(t+c),\qquad
 F_i=G_0(a_i-z_i)G_0(a_i+z_i),\qquad
 H_{a,i}=\tfrac12\partial F_i/\partial a_i,\qquad
 u_i=-12F_i/H_{a,i}.
\]
The derivative is taken with \(z_i\) formally held fixed,
\emph{before} substituting \(z_i=a_i(z+100i)\).
The balancing factors are
\begin{equation}\label{eq:abc-balance}
 R_{ir}=\frac{C_{ir}^2}{u_i(1+u_i^2C_{ir}^2)},\qquad
 W_i=\prod_{r=1}^{6}(1+R_{ir}^2).
\end{equation}

For completeness, a common bounded neighborhood with nonzero
tests can be chosen before selecting the slope. Choose a fixed
allowed \(U_0=[j_0]P\), with \(64\mid j_0\), away from every
forbidden point of \(U\) and \(U-A_i\), and with \(z(U_0)\ne0\)
at all the finitely many relevant embeddings. Choose \(H_0\)
sufficiently close to the identity at the prepared finite and
real places. The ratio \(z(U_0)/z(H_0)\) then has arbitrarily
large absolute value at each such place, whereas every ratio
\(\overline{z(U_0-A_i)/z(H_0-P)}\) stays bounded. The
denominator \(z(-P)\) is fixed and nonzero. The same ratio
also dominates the finitely many fixed \(n_i\). Thus all target
and paired-ratio tests are nonzero. The reciprocal-abscissa
tests are finite and nonzero after the point guards are avoided.
Take sufficiently small neighborhoods to preserve nonvanishing,
their finite valuations, and all the source chart conditions.
Include the finite restrictions in \(M\).

At every fixed prepared place and at real infinity,
\[
                    u_i\sim a_i(z+100i)^2 .
\]
For fixed nonzero \(C_{ir}\), \eqref{eq:abc-balance} therefore
tends to zero as the common slope tends to infinity, uniformly
on the chosen neighborhoods. An identically zero raw test,
if retained, contributes the factor one instead. In particular
one again does not require \(r\) to approach different \(n_i\)
at the prime two.

The source trace lemma gives
\[
 n_i^2T_i-57/4\in16\Z_2,\qquad
 v_2(T_i)=-2-2v_2(n_i),
\]
and its square-class conclusion makes \(a_i\) a
\(\Q_2\)-square. Hence the substituted contact parameter has
the same 2-adic square-class preparation as in the source.
All other effects of the finitely many fixed \(a_i\) are
absorbed in the fixed preparation and avoidance sets.
The prime-power CRT refinements again preserve all those
neighborhoods.

\section{Pole completion and both directions of the equivalence}
\label{sec:complete}

We now finish the ordered completeness proof. For integer inputs,
choose their private roots and the fixed common scalar/root
data as in Section~\ref{sec:fixed}. Choose safe initial \(H,U\)
neighborhoods and \(M\). The contact sets are now fixed.
Apply the retained prime supplier once to their union with
the avoidance and slope conditions stated above. Select
\(h,j\) by Lemma~\ref{lem:indices}, preserving the finite
preparations and real neighborhoods. Sections~\ref{sec:formal}
and \ref{sec:fixed} verify every changed local comparison.
All other source comparisons are unchanged.

For the general routes the actual pole inputs are
\[
             e_{i\ell}=
       \frac{4\rho_iR_{i\ell}}
            {\Tr(W_{i\ell}^{-1})^2}.
\]
The source parity argument gives square reduced denominator.
At two these lie in \(7\Q_2^{\times2}\) with depth at least
ten, and the real preparation gives
\[
                    0<e_{i\ell}<8\rho_i<1 .
\]
The new numerator preserves every input to that parity and
trace-unit argument. For the ordinary-\(abc\) route the
corresponding actual inputs are
\[
         e_i=\frac{28u_i}{(1+u_i)^2\Tr(W_i^{-1})^2},
                       \qquad d_i=7u_i .
\]
They have exactly the v88 square-denominator, class-seven,
depth, and \(0<e_i<1\) properties. None of these values is
identified between inputs.

For every detector retain the existing guarded square
extraction \(e=dk^2\) and the v88 normalized product equation
\[
                 Y^2=(1-e)(v^4-d^2)(s^4-d^2).
\]
Its two parameters \(v,s\) are private rational unknowns.
The v88 transfer from Bhargava--Burungale Lemma 2.2 and
Theorem 2.3 supplies them independently with a dense choice
of allowed parameter pairs. The finite residual pole guards
can therefore be avoided. The terminal root \(Y=0\) is
permitted; no terminal nonzero guard is introduced.
This proves simultaneous completeness at the level of the
guarded point and pole equations.

Conversely, project any rational solution of the new equations
to one original source block. Every rational expression,
point relation, detector relation, and denominator guard of
that block is still present after the common-variable
substitution. The original soundness theorem therefore
forces the block's input to be an integer. In the ordinary-\(abc\)
route it first forces \(X_i^2+1\) to be an integer, and the
elementary reduced-denominator argument in
Section~\ref{sec:roles} forces \(X_i\in\Z\).
This argument is valid for all rational solutions, not merely
the constructed integer witnesses. It proves that the joint
input projection is exactly \(\Z^k\).

\section{Finite algebra and complete polynomial specification}
\label{sec:polynomial}

We recall enough of the finite construction to specify the
output unambiguously. A root point has two retained rational
coordinates \(h_0,h_1\) and uses
\[
 h=h_0+h_1\sqrt2,\qquad x=h^2,\qquad
 j^2=h^4-D,\qquad y=hj.
\]
An auxiliary point instead retains the two rational components
of \(x\), with \(j^2=x^3-Dx\), \(y=j\).
In either case write the square input as \(A+B\sqrt2\).
Its free rank-four algebra has basis \(1,U,V,W\) and table
\[
\begin{array}{lll}
 U^2=2A+2W,& V^2=A-W,& UV=2B,\\
 UW=AU-2BV,& VW=BU-AV,& W^2=A^2-2B^2.
\end{array}
\]
The square root is \((U+V\sqrt2)/2\). This table is valid
at every base specialization, including nonreduced ones.
Each terminal pole supplies a monic quadratic stage of rank two.

Take one point algebra for each common role and one for each
private root \(A_i\), tensor them, and adjoin all the private
terminal quadratic stages. Let \(g\) be the product of every
retained source guard numerator after the common-numerator
substitution. The construction of \cite{shared} handles
guarded rational denominators without requiring them to be
nonzero at unselected characters. It gives an integer
polynomial \(\Phi_k\), using just one additional rational
encoder coordinate, such that
\[
 \exists U_*\in\Q\ \Phi_k(\mathbf X,\mathbf w,U_*)=0
 \quad\Longleftrightarrow\quad
 \exists\phi:A\longrightarrow\Q\quad\phi(g)\ne0 .
\]
No reducedness or monogenicity hypothesis is used. In particular,
any number of terminal zero roots is allowed.

For a directly usable specification, if \(A\) is the finite
algebra of rank \(n\), choose generating elements
\(\theta_1,\ldots,\theta_m\), put
\(\gamma_0=g,\gamma_r=g\theta_r\), and let \(A_r\) be their
multiplication matrices. Set \(\ell=m+1\), \(J=\ell(n+1)^2\),
and
\[
\begin{split}
 \widetilde A_r&=\operatorname{diag}(A_r,0),\qquad
 D_r=\widetilde A_r\otimes I-I\otimes\widetilde A_r,\\
 \prod_{r=0}^{m}\det(TI-D_r)
       &=T^J+a_1T^{J-1}+\cdots+a_J,\qquad a_0=1,\\
 R_0&=2+(n+1)+\sum_{r,v,w}(A_r)_{vw}^2+\sum_{j=1}^{J}a_j^2,\\
 B_0&=(n+1)R_0,\qquad K_0=R_0(2B_0)^{J-1}.
\end{split}
\]
For \(0\le s\le J\) put
\[
 \Delta_s=a_s^2,\quad
 C_s=\Delta_s+(4\ell B_0+2)K_0,\quad
 M_s=\Delta_s\sum_{r=0}^{m}C_s^{\,r+1}\Delta_s^{\,m-r}A_r .
\]
The required encoder is
\begin{equation}\label{eq:encoder}
 \Phi_k=
 \prod_{s=0}^{J}
 \left(\det(I-U_*M_s)^2+\sum_{j=s+1}^{J}a_j^2\right).
\end{equation}
After the guarded denominator normalization, these are
polynomials with integer coefficients. The proof in
\cite{shared} separates the values of all permitted geometric
characters from each other and from zero on the actual last
nonzero coefficient stratum. A rational determinant root then
selects a character fixed by Galois, hence a rational character.
Nilpotent multiplicities do not affect the argument.

\section{Coordinate and degree ledgers}
\label{sec:counts}

For the general routes the retained rational coordinates are
\[
\begin{array}{lr}
 \text{inputs}&k\\
 \text{private root charts }A_i&2k\\
 \text{eight common point charts}&16\\
 \text{one common contact parameter}&1\\
 \text{private product-pole parameters}&2ek\\
 \text{one finite-algebra encoder}&1.
\end{array}
\]
Thus the total is \((3+2e)k+18\). The finite algebra has rank
\[
                   4^{k+8}2^{ek}=2^{16+(2+e)k}.
\]
For ordinary \(abc\) there are two common point charts, so
the total is
\[
                    k+2k+4+1+2k+1=5k+6,
\]
and the finite rank is \(4^{k+2}2^k=2^{4+3k}\).

We use the mixed-integer universal polynomial with eight
integer unknowns and one nonnegative integer unknown
\cite[Theorem 1.1(i)]{SunHTP}. The latter is represented by
\[
                         L_+(s,t,u)=s^2+t^2+u^2+u.
\]
Every integral rational value of \(L_+\) is nonnegative since
\(L_+\ge-1/4\), and every \(n\ge0\) occurs: represent
\(4n+1\) by three integer squares, exactly one odd, and
divide the resulting identity by four. Substituting \(L_+\)
for the ninth input in both the universal polynomial and
\(\Phi_9\) replaces one coordinate by three. Combining the
two equalities by a sum of squares adds no coordinate.
The totals are therefore
\[
                 (7\cdot9+18)+2=83,\qquad
                 (9\cdot9+18)+2=101,\qquad
                 (5\cdot9+6)+2=53.
\]
The instance integer occurs in the coefficients, not among
the unknowns.

For completeness the operator-degree induction from
\cite{shared} gives the following exact ledger.
Assign weight one to retained coordinates, weight zero to
fixed coefficient parameters, maximum at an addition, and
sum at a product. The point table gives multiplication
degree bounds \(a,a,2a\) for \(U,V,W\) when
\(\deg A,\deg B\le a\). A terminal monic equation
\(q^2+c=0\) gives degree bound \(\deg M_c\) for
multiplication by \(q\). Its coefficients and every guard
are checked not to depend on terminal roots. Tensor
extensions leave preceding bounds unchanged.

The point algebra over \(\Q\) is generated by \(U,V\), since
\(W=U^2/2-A\); these two generators and every terminal root
suffice for \eqref{eq:encoder}. If \(d_g\) is the sum of the
guard operator-degree bounds and \(d_\theta\) is the largest
generator bound, one may take \(d=d_g+d_\theta\).
\begin{center}
\begin{tabular}{lrrrrr}
\toprule
Route&rank \(n\)&\(\ell\)&\(d_g\)&\(d_\theta\)&\(d\)\\
\midrule
General&\(2^{52}\)&53&2201826&100232&2302058\\
Specified&\(2^{61}\)&62&4025982&100232&4126214\\
Ordinary \(abc\)&\(2^{31}\)&32&3103028&206888&3309916\\
\bottomrule
\end{tabular}
\end{center}
The joint-certificate degree is bounded by
\begin{equation}\label{eq:degree}
 D_{\Phi}=2n(J+1)\left[1+d\{1+2J(1+\ell(J+1))\}\right],
             \qquad J=\ell(n+1)^2 .
\end{equation}
For a fully bounded mixed universal source use the proved
universal pair \((58,4)\) in
\cite[Theorem 3, Definition 5.3 and Section 6.1]{BDHMS}.
After the nonnegative wrapper its degree is at most
\[
 \eta(\nu,\delta)=15616+233856\delta+
       (233952\delta+467712\delta^2)(2\delta+1)^{\nu+1}.
\]
The final sum of squares thus has degree at most
\(\max\{2\eta(58,4),4D_\Phi\}\). Substituting the exact
integers above gives
\[
\begin{array}{c|c|c}
 \text{route}&\text{rational unknowns}&\text{degree strictly below}\\ \hline
 \text{general}&83&2^{413}\\
 \text{specified}&101&2^{477}\\
 \text{ordinary }abc&53&2^{263}.
\end{array}
\]
These are certified upper bounds, not attained degrees or
claims of optimality.

\section{Reproducibility and precise scope}

The complete new specifications are in
\nolinkurl{undecidability/common_numerator_sources/}, with
one \texttt{common\_source.json.gz} and one certificate for
each of \texttt{general25}, \texttt{explicit27}, and
\texttt{abc11}. The builder
\nolinkurl{undecidability/code/build_common_numerator.py}
starts from the published joint coefficient circuits.
It identifies exactly the two retained coordinates and
three finite generators named \texttt{blocki.UA} with
\texttt{shared.UA}. It substitutes through the entire
coefficient circuit, unifies only identical point blocks,
deduplicates identical guards, and composes every original
source-coordinate and coefficient-node map.

The independent replay
\nolinkurl{undecidability/verify_common_numerator.py}
checks the complete transformation, source projections,
point multiplication identities, count and dependency
conditions, and the integer degree bounds. Its certificate is
\nolinkurl{undecidability/certificates/common_numerator_transfers.json}.
The numbers of coefficient nodes are respectively
41491, 58303, and 12890; the numbers of distinct retained
guard nodes are 635, 905, and 206. The retained tuple
coordinate counts, before the mixed wrapper, are 81, 99,
and 51. These are finite, exact specifications; the final
determinant polynomials have not been expanded.

Changing an ineffective fixed parameter to an actual
defining member requires rebuilding the corresponding
v88 native source and applying the same two sharing
transformations. The circuit degrees are unchanged:
\(L\) enters coefficient constants such as \(67^L\), not
exponents of rational coordinates, and \(B\) is a fixed
coefficient parameter. Therefore the stated degree
ceilings hold for actual defining members as well as
for the algebraic fixtures. Finite verification does not
certify the arithmetic projection of a fixture.

Finally, the new congruence argument applies to numerator
roles belonging to distinct input contact families. It
does not identify the different label numerators inside
one input block: a single opposite annihilator would then
receive potentially different labels modulo the same
modulus. Nor does this construction change the quartic
degree restriction. It supplies a proved reduction in the
number of variables for general fixed degree under the
retained source premises, and makes no global minimality
claim.

\begin{thebibliography}{9}
\bibitem{shared}
\(\mathbb Z\mathbb Q\) Diophantine Project,
\emph{Shared integer certificates and rational undecidability in
99, 117, and 69 variables}, 23 November 2026.
\nolinkurl{undecidability/papers/shared_integer_transfers.tex}.

\bibitem{v78}
\(\mathbb Z\mathbb Q\) Diophantine Project,
\emph{Four-contact scalar reduction}, v78, and
\emph{Twelve-contact first-scale construction}, v78 (2026).
\nolinkurl{papers/few_variables25/four_contact_scalar_20261119_v78.tex};
\nolinkurl{papers/conditional_few_variables/twelve_contact_first_scale_20261119_v78.tex}.

\bibitem{v88}
\(\mathbb Z\mathbb Q\) Diophantine Project,
\emph{Product poles and lower degrees in the few-variable
constructions}, v88 (2026).
\nolinkurl{papers/few_variables25/bb_product_degree_transfer_20261123_v88.tex}.

\bibitem{patterns}
Pinned Sun source, \texttt{sun-2610.20691v2-P103q.tex},
Theorem 4.1, label \texttt{an:thm:patterns}, upstream commit
\texttt{14e8f601fe0d9173d4a5f084aedc3e07066a7688}.
\url{https://github.com/JoseBrox/diophantine-z-q_improvements/blob/14e8f601fe0d9173d4a5f084aedc3e07066a7688/quartics/quartic-1272/sources/sun-2610.20691v2-P103q.tex}.
The read source has SHA256
\nolinkurl{b561d39ffc0e46798bd6ad5d3cfb8b543b75871780758876a8af67fb6daf15ba}.
Its all-sign order assertion is already part of the prime supplier
cited in the project arithmetic dependencies.

\bibitem{BB}
M. Bhargava and A. Burungale, the square-denominator and
product-pole results audited in v88,
\url{https://arxiv.org/abs/2611.16130v2},
Lemma 2.2 and Theorem 2.3.

\bibitem{SunHTP}
Z.-W. Sun, \emph{Further results on Hilbert's Tenth Problem},
Science China Mathematics \textbf{64} (2021), 281--306,
\url{https://arxiv.org/abs/1704.03504}, Theorem 1.1(i).

\bibitem{BDHMS}
J. Bayer, M. David, M. Ha{\ss}ler, Y. Matiyasevich and
D. Schleicher,
\emph{Diophantine equations over the integers: universal
bounds and parallel formalization}, version 3, 22 October 2026,
\url{https://arxiv.org/abs/2506.20909v3}.
\end{thebibliography}
\end{document}
